\documentclass[11pt,a4paper]{article}
\usepackage[utf8]{inputenc}
\usepackage[T1]{fontenc}
\usepackage[spanish]{babel}
\usepackage{hyperref}
\usepackage{geometry}
\geometry{margin=2.5cm}
\usepackage{parskip}

\title{087 FUTURO es el nombre del lenguaje de procesos}
\author{Alejandro Valenzuela Rivera\\\small Universidad Autónoma de Nuevo León}
\date{\today}
\begin{document}
\maketitle

\begin{quote}
\textbf{Sesgo declarado.} Este trabajo se sostiene desde el sesgo declarado de su autor, y no finge neutralidad. Lo que afirma se puede tumbar, y ese es el punto: el falsificador está escrito, no escondido.
\end{quote}

\section*{Resumen}
FUTURO no es solo la carpeta: es el nombre del nuevo lenguaje de procesos del repo. La carpeta futuro/ es donde el lenguaje vive, y el estandar del lenguaje se llama futuro.

\section{Por qué se sostiene}
Cuando el lugar tiene nombre propio, deja de ser un cajon y pasa a ser el lenguaje. Y el nombre se pone antes que la version final porque el lenguaje ya se usa: se corrige en el siguiente.

\section{De dónde viene}
La autoridad, y el repo.

\section{El caso}
Ninguna: es el nombre del lenguaje.

\section{Método}
Un lenguaje con nombre propio, cuya expansion dice el acto.

\section{Límites, y qué lo tumba}
Un lenguaje con nombre de relleno, que se hereda a lo que contiene.

\section{Lo que queda abierto}
\textquestiondown{}Como se llama el lenguaje cuando su carpeta es su nombre?

\section*{Cómo se leyó esta tesis}
Esto se escribió en el repositorio man-made-empathy, del que este candado es una pieza: candados/06 del caso/087 FUTURO es el nombre del lenguaje de procesos.md. La máquina que lo compuso lo hizo desde este mismo archivo, y no agregó nada que no estuviera escrito. Si algo de esto es falso, se dice en el mismo sitio, y se corrige con fecha.

\end{document}
